\chapter{Migraine}

\section*{Definition and Overview}
Migraine is a common neurological disorder characterised by recurrent episodes of moderate to severe headache, typically unilateral and pulsating, often accompanied by nausea, vomiting, and sensitivity to light, sound, or smell. It is a complex brain network disorder involving the trigeminovascular system, cortical spreading depression, and altered sensory processing. Migraine affects approximately 15\% of the global population, with a 3:1 female-to-male predominance. Attacks last 4--72 hours if untreated and are a leading cause of disability in people under 50. Migraine can be episodic (fewer than 15 headache days per month) or chronic (15 or more headache days per month for over 3 months, with at least 8 having migraine features).

\section*{Epidemiology}
Migraine affects around 4.9 million Australians, with peak prevalence between ages 25 and 55. It is three times more common in women than men, largely due to hormonal influences. Chronic migraine affects 1--2\% of the population. The World Health Organization ranks migraine as the second leading cause of years lived with disability globally. Migraine has a strong genetic component; having a first-degree relative with migraine increases risk 2--3 fold. Comorbidities include depression, anxiety, insomnia, and cardiovascular disease. The economic burden in Australia exceeds \$35 billion annually in lost productivity and healthcare costs.

\section*{Aetiology and Pathophysiology}
\begin{itemize}
    \item \textbf{Genetics:} polygenic inheritance with multiple susceptibility loci (e.g., \textit{CACNA1A}, \textit{ATP1A2}, \textit{SCN1A} for familial hemiplegic migraine); common variants identified through GWAS
    \item \textbf{Cortical spreading depression (CSD):} wave of neuronal and glial depolarisation followed by suppression, underlying aura and activating trigeminal afferents
    \item \textbf{Trigeminovascular activation:} CSD and other triggers release CGRP, substance P, and PACAP from trigeminal nerve endings, causing neurogenic inflammation and vasodilation
    \item \textbf{CGRP pathway:} calcitonin gene-related peptide is the key neuropeptide; elevated during attacks and central to modern preventive therapies
    \item \textbf{Brainstem and hypothalamic involvement:} modulation of pain processing, autonomic symptoms, and premonitory features (yawning, cravings, mood changes)
    \item \textbf{Triggers:} hormonal fluctuation, stress, sleep disruption, fasting, alcohol (especially red wine), aged cheese, caffeine withdrawal, weather changes, strong odours, bright lights
\end{itemize}

\section*{Clinical Features}
\begin{itemize}
    \item \textbf{Premonitory phase (hours to days before):} fatigue, mood change, neck stiffness, food cravings, yawning, urinary frequency
    \item \textbf{Aura (20--60 min, in 25--30\%):} visual (scintillating scotoma, fortification spectra), sensory (tingling spreading over face/arm), speech/language disturbance, motor weakness (hemiplegic migraine)
    \item \textbf{Headache phase:} unilateral or bilateral pulsating pain, moderate to severe intensity, aggravated by routine physical activity; accompanied by nausea/vomiting, photophobia, phonophobia, osmophobia
    \item \textbf{Postdrome (``migraine hangover''):} fatigue, cognitive difficulty, mood change, residual neck stiffness lasting hours to a day
    \item \textbf{Chronic migraine:} 15+ headache days/month, often with daily background headache and superimposed migraine attacks; frequently associated with medication overuse
\end{itemize}

\section*{Differential Diagnosis}
\begin{itemize}
    \item \textbf{Tension-type headache:} bilateral, pressing/tightening, mild-moderate, no nausea, no aggravation by activity
    \item \textbf{Cluster headache:} strictly unilateral, severe orbital/supraorbital/temporal, 15--180 min, with ipsilateral autonomic features (lacrimation, nasal congestion, ptosis); circadian periodicity
    \item \textbf{Medication-overuse headache:} daily/near-daily headache in setting of frequent acute medication use ($\ge$10--15 days/month); improves after withdrawal
    \item \textbf{Sinus headache:} often misdiagnosed; true sinusitis has purulent discharge, fever, and responds to antibiotics
    \item \textbf{Secondary headaches:} thunderclap (subarachnoid haemorrhage), progressive with neurological signs (tumour), postural (low CSF pressure), with fever/stiff neck (meningitis)
    \item \textbf{New daily persistent headache (NDPH):} abrupt onset, unremitting daily headache without migraine features
\end{itemize}

\section*{When to Seek Medical Care}
See a doctor for frequent headaches (more than 4 days/month), headaches that interfere with daily activities, or a change in headache pattern. First-line preventive therapy should be offered for 4+ migraine days/month or disabling attacks.

\textbf{Call 000 (or local emergency services) for:}
\begin{itemize}
    \item Sudden, severe ``thunderclap'' headache (worst ever) -- possible subarachnoid haemorrhage
    \item Headache with fever, stiff neck, altered consciousness -- possible meningitis
    \item Headache with focal weakness, speech difficulty, vision loss -- possible stroke
    \item Headache after head trauma with confusion, vomiting, or drowsiness
    \item New headache in pregnancy or after age 50
    \item Headache precipitated by Valsalva, exertion, or sexual activity
\end{itemize}

\section*{Diagnosis and Investigations}
\begin{itemize}
    \item \textbf{Clinical diagnosis:} based on ICHD-3 criteria; no biomarker or imaging test confirms migraine
    \item \textbf{Headache diary:} essential for diagnosis, trigger identification, and treatment monitoring (frequency, severity, duration, medication use, menstrual cycle)
    \item \textbf{Neurological examination:} normal between attacks; may show transient deficits during hemiplegic migraine aura
    \item \textbf{Imaging:} MRI brain indicated for atypical features, abnormal exam, new headache after 50, thunderclap onset, or progressive pattern; CT if MRI unavailable
    \item \textbf{Blood tests:} not routine; consider ESR/CRP for giant cell arteritis in patients $>$50 with new headache, or thyroid function if clinically indicated
    \item \textbf{Referral:} neurologist or headache specialist for diagnostic uncertainty, treatment failure, chronic migraine, hemiplegic migraine, or need for advanced therapies
\end{itemize}

\section*{Management}
\begin{itemize}
    \item \textbf{Acute treatment -- early is essential:}
    \begin{itemize}
        \item Simple analgesics: aspirin 900 mg or ibuprofen 400--600 mg or paracetamol 1000 mg, preferably soluble/liquid formulations
        \item Anti-emetics: metoclopramide 10 mg or prochlorperazine 10 mg (also enhance gastric absorption)
        \item Triptans (5-HT$_{1B/1D}$ agonists): sumatriptan, rizatriptan, eletriptan, zolmitriptan; first-line for moderate-severe attacks; take at onset of pain; avoid in cardiovascular disease
        \item Ditans (lasmiditan): 5-HT$_{1F}$ agonist for patients with cardiovascular contraindications to triptans
        \item Gepants (rimegepant, ubrogepant): oral CGRP receptor antagonists; no cardiovascular restrictions; useful when triptans contraindicated or ineffective
        \item Avoid opioids and barbiturates (high risk of medication-overuse headache and dependence)
    \end{itemize}
    \item \textbf{Preventive treatment (offered for $\ge$4 migraine days/month or disabling attacks):}
    \begin{itemize}
        \item First-line oral: beta-blockers (propranolol, metoprolol), topiramate, amitriptyline, candesartan
        \item Monoclonal antibodies against CGRP or its receptor: erenumab, fremanezumab, galcanezumab, eptinezumab (subcutaneous/IV quarterly); high efficacy, good tolerability, PBS-listed in Australia for chronic migraine failing 3 preventives
        \item Oral gepants (atogepant, rimegepant): daily or alternate-day for prevention
        \item OnabotulinumtoxinA: 155--195 units every 12 weeks for chronic migraine (PBS-listed)
        \item Neuromodulation devices: non-invasive vagus nerve stimulation, supraorbital nerve stimulation
    \end{itemize}
    \item \textbf{Medication-overuse headache:} withdrawal of overused medication (bridge with prednisolone or IV dihydroergotamine if needed), then start preventive
    \item \textbf{Lifestyle:} regular sleep, meals, hydration; aerobic exercise; stress management (CBT, mindfulness, biofeedback); trigger avoidance
    \item \textbf{Menstrual migraine:} perimenstrual prevention with frovatriptan, naratriptan, or magnesium; or continuous hormonal contraception to suppress cycle
    \item \textbf{Pregnancy:} avoid most preventives; paracetamol safe; metoclopramide/prochlorperazine; some triptans (sumatriptan) have reassuring registry data; discuss with specialist
\end{itemize}

\section*{Complications}
\begin{itemize}
    \item \textbf{Chronic migraine:} progression from episodic to chronic (2.5\% per year); associated with depression, anxiety, and medication overuse
    \item \textbf{Medication-overuse headache:} develops with $\ge$10 days/month of triptans/ergots/opioids/combination analgesics or $\ge$15 days/month of simple analgesics
    \item \textbf{Status migrainosus:} debilitating attack $>$72 hours; may require IV therapy
    \item \textbf{Migrainous infarction:} ischaemic stroke during typical aura; rare ($<$1 per 100,000 person-years)
    \item \textbf{Persistent aura without infarction:} aura lasting $>$1 week without stroke on imaging
    \item \textbf{Psychiatric comorbidity:} 2--5x increased risk of depression, anxiety, bipolar disorder
    \item \textbf{Cardiovascular risk:} migraine with aura associated with 2x relative risk of ischaemic stroke; combined oral contraceptive further increases risk
\end{itemize}

\section*{Prognosis and Outcomes}
Migraine is a chronic condition with a variable course. Many patients improve with age, particularly after menopause in women. With appropriate acute and preventive treatment, most patients achieve significant reduction in attack frequency and disability. The introduction of CGRP-targeted therapies has transformed outcomes for refractory patients. Early diagnosis, patient education, avoidance of medication overuse, and a structured management plan are key to optimal long-term outcomes.

\section*{Prevention and Self-Care}
\begin{itemize}
    \item Keep a headache diary (paper or app) to identify triggers and track treatment response
    \item Maintain regular sleep-wake cycle, even on weekends
    \item Eat regular meals; avoid prolonged fasting
    \item Stay well hydrated (2--3 L water daily)
    \item Limit caffeine to $<$200 mg/day and avoid sudden withdrawal
    \item Regular aerobic exercise (150 min/week moderate intensity)
    \item Stress management: mindfulness, CBT, relaxation training, biofeedback
    \item Take acute medication at first sign of pain, but limit to $<$10 days/month to avoid overuse
    \item Consider preventive therapy if 4+ migraine days/month or attacks are disabling
    \item Review medications annually; step down preventives after 6--12 months of good control
    \item Discuss contraception and pregnancy plans with GP/neurologist
    \item Join support organisations: Migraine \& Headache Australia, Headache Australia
\end{itemize}

\section*{Sources and Further Reading}
The following reputable sources provide further information for healthcare students and patients seeking reliable, current advice about migraine.

\begin{itemize}
    \item Migraine \& Headache Australia. \textit{Migraine information and support.} https://www.migraine.org.au/
    \item Headache Australia (Brain Foundation). \textit{Migraine and headache resources.} https://headacheaustralia.org.au/
    \item International Headache Society. \textit{ICHD-3 Classification and guidelines.} https://ihs-headache.org/
    \item Australian and New Zealand Association of Neurologists. \textit{Migraine management position statements.} https://www.anzan.org.au/
    \item NPS MedicineWise. \textit{Migraine treatment options.} https://www.nps.org.au/
    \item The Migraine Trust (UK). \textit{Patient information and research.} https://www.migrainetrust.org/
    \item American Headache Society. \textit{Clinical practice guidelines.} https://americanheadachesociety.org/
\end{itemize}